\documentclass[12pt,a4paper]{article}
\usepackage[utf8]{inputenc}
\usepackage[T1]{fontenc}
\usepackage[french]{babel}
\usepackage{geometry}
\usepackage{listings}
\usepackage{xcolor}
\usepackage{booktabs}
\usepackage{titlesec}
\usepackage{hyperref}
\usepackage{parskip}
\usepackage{fancyhdr}

\geometry{margin=2.5cm}

% Colors
\definecolor{graphqlpink}{RGB}{225,0,152}
\definecolor{codebg}{RGB}{245,245,245}
\definecolor{keywordcolor}{RGB}{0,0,180}
\definecolor{stringcolor}{RGB}{163,21,21}
\definecolor{commentcolor}{RGB}{0,128,0}

% GraphQL listing style
\lstdefinelanguage{GraphQL}{
  keywords={type, input, enum, interface, union, scalar, query, mutation, subscription, fragment, on, true, false, null},
  keywordstyle=\color{keywordcolor}\bfseries,
  stringstyle=\color{stringcolor},
  commentstyle=\color{commentcolor}\itshape,
  morestring=[b]",
  morecomment=[l]{\#},
  sensitive=true
}

% TypeScript listing style
\lstdefinelanguage{TypeScript}{
  keywords={const, let, var, function, return, import, export, from, interface, type, extends, implements, class, new, async, await, true, false, null, undefined},
  keywordstyle=\color{keywordcolor}\bfseries,
  stringstyle=\color{stringcolor},
  commentstyle=\color{commentcolor}\itshape,
  morestring=[b]",
  morestring=[b]',
  morestring=[b]`,
  morecomment=[l]{//},
  morecomment=[s]{/*}{*/},
  sensitive=true
}

\lstset{
  backgroundcolor=\color{codebg},
  basicstyle=\ttfamily\small,
  breaklines=true,
  frame=single,
  rulecolor=\color{gray!40},
  numbers=left,
  numberstyle=\tiny\color{gray},
  numbersep=8pt,
  tabsize=2,
  showstringspaces=false,
  xleftmargin=15pt
}

% Title formatting
\titleformat{\section}{\large\bfseries\color{graphqlpink}}{}{0em}{}[\titlerule]
\titleformat{\subsection}{\normalsize\bfseries}{}{0em}{}

% Header/Footer
\pagestyle{fancy}
\fancyhf{}
\rhead{\small Rapport Technique --- TP2 GraphQL}
\lhead{\small Personne 1 --- Schéma GraphQL}
\rfoot{\thepage}

\hypersetup{
  colorlinks=true,
  linkcolor=graphqlpink,
  urlcolor=graphqlpink
}

% ─────────────────────────────────────────────
\begin{document}

\begin{titlepage}
  \centering
  \vspace*{3cm}
  {\Huge\bfseries Rapport Technique\\[0.4em]
  \color{graphqlpink} TP2 --- GraphQL\par}
  \vspace{1.5cm}
  {\large Personne 1 : Définition du Schéma GraphQL\\
  Types, Relations \& Enum\par}
  \vspace{1cm}
  \rule{0.6\textwidth}{0.4pt}\\[1em]
  {\normalsize
    Module : Développement Web\\
    Encadrant : Aymen Sellaouti\\
    Date : \today
  }
  \vfill
\end{titlepage}

\tableofcontents
\newpage

% ─────────────────────────────────────────────
\section{Stack Technique Choisie}

\begin{center}
\begin{tabular}{@{}lp{10cm}@{}}
\toprule
\textbf{Outil} & \textbf{Rôle} \\
\midrule
TypeScript    & Typage statique, évite les erreurs au runtime \\
GraphQL Yoga  & Serveur GraphQL moderne, plus léger qu'Apollo, supporte nativement les Subscriptions via SSE \\
Node.js HTTP  & Couche transport sous GraphQL Yoga \\
Nodemon + ts-node & Rechargement automatique en développement \\
\bottomrule
\end{tabular}
\end{center}

% ─────────────────────────────────────────────
\section{Architecture du Projet}

\begin{lstlisting}[language=bash, numbers=none, caption={Structure des dossiers}]
src/
├── schema/
│   └── typeDefs.ts               <- Personne 1 (schéma GraphQL)
├── data/
│   └── db.ts                     <- Personne 2 (base de données fictive)
├── context/
│   └── context.ts                <- Personne 2 (configuration du context)
├── resolvers/
│   ├── index.ts                  <- point de fusion des resolvers
│   ├── query/cvQueries.ts        <- Personne 3 (queries)
│   ├── types/cvTypeResolvers.ts  <- Personne 4 (field resolvers)
│   ├── mutation/cvMutations.ts   <- Personne 5 (mutations)
│   └── subscription/cvSubscriptions.ts <- Personne 5 (subscriptions)
└── server.ts                     <- point d'entrée du serveur
\end{lstlisting}

Le principe appliqué est la \textbf{séparation des responsabilités} : chaque fichier a un rôle unique et précis. C'est l'application du design pattern \textbf{Single Responsibility Principle}.

% ─────────────────────────────────────────────
\section{Fichier Principal --- \texttt{src/schema/typeDefs.ts}}

C'est le cœur du travail de Personne 1. En GraphQL, les \texttt{typeDefs} définissent le \textbf{contrat de l'API} --- ce que le client peut demander et ce que le serveur s'engage à retourner. Tout est écrit en \textbf{SDL (Schema Definition Language)}.

\subsection{3.1 --- \texttt{enum Role}}

\begin{lstlisting}[language=GraphQL, caption={Définition de l'enum Role}]
enum Role {
  USER
  ADMIN
}
\end{lstlisting}

Un \texttt{enum} est un type dont la valeur est restreinte à une liste fixe. Ici \texttt{Role} ne peut valoir que \texttt{USER} ou \texttt{ADMIN}, jamais une autre chaîne de caractères. Le professeur l'a explicitement demandé dans le sujet (page 3). Sans enum, on utiliserait un \texttt{String} libre ce qui est moins sûr et moins expressif.

\subsection{3.2 --- \texttt{type User}}

\begin{lstlisting}[language=GraphQL, caption={Type User}]
type User {
  id: ID!
  name: String!
  email: String!
  role: Role!
}
\end{lstlisting}

\begin{itemize}
  \item \texttt{ID!} --- identifiant unique. Le \texttt{!} signifie \textbf{non-nullable} : toujours présent, jamais null.
  \item \texttt{String!} --- chaîne de caractères obligatoire.
  \item \texttt{role: Role!} --- utilise l'enum défini ci-dessus, obligatoire.
\end{itemize}

\subsection{3.3 --- \texttt{type Skill}}

\begin{lstlisting}[language=GraphQL, caption={Type Skill}]
type Skill {
  id: ID!
  designation: String!
}
\end{lstlisting}

Type simple avec deux champs. \texttt{designation} correspond au nom du skill (ex: ``React'', ``Java'', ``Spring Boot'').

\subsection{3.4 --- \texttt{type Cv} --- le type central avec ses relations}

\begin{lstlisting}[language=GraphQL, caption={Type Cv avec relations}]
type Cv {
  id: ID!
  name: String!
  age: Int!
  job: String!
  user: User
  skills: [Skill!]!
}
\end{lstlisting}

C'est ici que les \textbf{relations du schéma UML sont traduites en GraphQL} :

\begin{itemize}
  \item \texttt{user: User} --- relation \textbf{many-to-one} : un CV appartient à un User. Nullable (pas de \texttt{!}) car un CV pourrait exister sans user assigné.
  \item \texttt{skills: [Skill!]!} --- relation \textbf{many-to-many} : un CV possède plusieurs Skills. La double notation signifie : la liste elle-même est non-nullable (\texttt{!} extérieur) et chaque élément de la liste est également non-nullable (\texttt{!} intérieur dans \texttt{[Skill!]}).
\end{itemize}

\textbf{Point important :} dans la base de données (\texttt{db.ts}), le CV stockera \texttt{userId} et \texttt{skillIds} comme de simples IDs. C'est Personne 4 qui écrira les resolvers qui transformeront ces IDs en vrais objets \texttt{User} et \texttt{Skill} lors de la requête. Personne 1 définit uniquement le \textbf{contrat GraphQL visible par le client}.

\subsection{3.5 --- \texttt{type Query}}

\begin{lstlisting}[language=GraphQL, caption={Déclaration des Queries}]
type Query {
  cvs: [Cv!]!
  cv(id: ID!): Cv
}
\end{lstlisting}

\begin{itemize}
  \item \texttt{cvs} --- retourne tous les CVs. La liste ne peut pas être null mais peut être vide.
  \item \texttt{cv(id: ID!)} --- retourne un CV par son ID. Nullable car si l'ID n'existe pas, on retourne \texttt{null}.
\end{itemize}

Ces signatures définissent le contrat, mais les \textbf{implémentations} (la logique métier) sont réalisées par Personne 3 dans \texttt{cvQueries.ts}.

\subsection{3.6 --- \texttt{input CreateCvInput} et \texttt{input UpdateCvInput}}

\begin{lstlisting}[language=GraphQL, caption={Inputs pour les mutations}]
input CreateCvInput {
  name: String!
  age: Int!
  job: String!
  userId: ID!
  skillIds: [ID!]!
}

input UpdateCvInput {
  name: String
  age: Int
  job: String
  userId: ID
  skillIds: [ID!]
}
\end{lstlisting}

Un \texttt{input} est un type spécial utilisé uniquement pour \textbf{envoyer des données} au serveur (dans les mutations). Il ne peut pas contenir de fields résolus.

Différence cruciale entre les deux :
\begin{itemize}
  \item \texttt{CreateCvInput} : tous les champs sont obligatoires (\texttt{!}) car pour créer un CV on a besoin de toutes les informations.
  \item \texttt{UpdateCvInput} : tous les champs sont optionnels (sans \texttt{!}) car lors d'une mise à jour on ne modifie que ce que l'on souhaite.
\end{itemize}

Les IDs des relations (\texttt{userId}, \texttt{skillIds}) sont passés comme données brutes --- cohérent avec la base relationnelle non imbriquée demandée par le professeur.

\subsection{3.7 --- \texttt{type Mutation}}

\begin{lstlisting}[language=GraphQL, caption={Déclaration des Mutations}]
type Mutation {
  createCv(input: CreateCvInput!): Cv!
  updateCv(id: ID!, input: UpdateCvInput!): Cv!
  deleteCv(id: ID!): Boolean!
}
\end{lstlisting}

\begin{itemize}
  \item \texttt{createCv} --- prend un \texttt{CreateCvInput} et retourne le CV créé.
  \item \texttt{updateCv} --- prend l'ID du CV à modifier ainsi que les nouvelles données, retourne le CV mis à jour.
  \item \texttt{deleteCv} --- prend un ID, retourne un \texttt{Boolean} (\texttt{true} si supprimé, \texttt{false} sinon).
\end{itemize}

Ces stubs permettent au projet de compiler immédiatement. Personne 5 n'aura qu'à écrire les resolvers correspondants.

\subsection{3.8 --- \texttt{type Subscription}}

\begin{lstlisting}[language=GraphQL, caption={Déclaration de la Subscription}]
type Subscription {
  cvChanged: Cv!
}
\end{lstlisting}

Une seule subscription \texttt{cvChanged} qui notifie à chaque ajout, modification ou suppression d'un CV. C'est le choix \textbf{économique en ressources} demandé page 6 du sujet : une subscription unique au lieu de trois (\texttt{cvAdded}, \texttt{cvUpdated}, \texttt{cvDeleted}) = moins de connexions ouvertes côté serveur. GraphQL Yoga utilise \textbf{SSE (Server-Sent Events)} au lieu de WebSockets, ce qui consomme encore moins de ressources.

% ─────────────────────────────────────────────
\section{Les Autres Fichiers (Structure)}

\subsection{4.1 --- \texttt{src/server.ts}}

\begin{lstlisting}[language=TypeScript, caption={Point d'entrée du serveur}]
const yoga = createYoga({
  schema: { typeDefs, resolvers } as any,
  context: createContext,
});
const server = createServer(yoga);
server.listen(4000, () => {
  console.log('Server running at http://localhost:4000/graphql');
});
\end{lstlisting}

GraphQL Yoga reçoit les \texttt{typeDefs} (travail de Personne 1), les \texttt{resolvers} (travail des autres membres) et le \texttt{context}. Il expose automatiquement une interface \textbf{GraphiQL} (interface de test interactive) sur \texttt{http://localhost:4000/graphql}.

\subsection{4.2 --- \texttt{src/context/context.ts}}

\begin{lstlisting}[language=TypeScript, caption={Configuration du Context}]
export const createContext = (): AppContext => ({
  db: { users, cvs, skills, cvSkills },
});
\end{lstlisting}

Le context est un objet \textbf{injecté automatiquement dans chaque resolver}. Grâce à lui, tous les resolvers accèdent à la base de données sans import direct --- c'est le pattern recommandé en GraphQL pour éviter le couplage fort entre les resolvers et la source de données.

\subsection{4.3 --- \texttt{src/resolvers/index.ts}}

\begin{lstlisting}[language=TypeScript, caption={Fusion des resolvers}]
export const resolvers = {
  Query:        { ...cvQueries },
  Mutation:     { ...cvMutations },
  Subscription: { ...cvSubscriptions },
  Cv:           { ...CvTypeResolvers },
};
\end{lstlisting}

Fusionne tous les resolvers via le \textbf{spread operator}. Chaque membre de l'équipe ajoute ses resolvers dans son propre fichier, et \texttt{index.ts} les assemble sans conflit --- c'est le pattern \textbf{Resolver Map}.

% ─────────────────────────────────────────────
\section{Résumé des Décisions Techniques}

\begin{center}
\begin{tabular}{@{}lp{9cm}@{}}
\toprule
\textbf{Élément} & \textbf{Décision technique} \\
\midrule
\texttt{enum Role}       & Restreint les valeurs possibles, plus sûr qu'un \texttt{String} libre \\
Relations dans \texttt{Cv} & Définies en GraphQL (\texttt{user}, \texttt{skills}), séparées des données brutes en DB \\
\texttt{[Skill!]!}       & Double non-nullable : ni la liste ni ses éléments ne peuvent être null \\
\texttt{input} séparés   & \texttt{Create} vs \texttt{Update} avec champs requis/optionnels selon le contexte \\
\texttt{cvChanged} unique & Une seule subscription pour 3 événements = économie de ressources \\
SSE via GraphQL Yoga     & Plus léger que WebSockets, répond à l'exigence page 6 du sujet \\
\bottomrule
\end{tabular}
\end{center}

\end{document}
